\section{Action, gravité et information thermique}
\label{sec:gravity}
\subsection{Période dodécaphasique et réalisation circulaire de l'action}
La longueur de la section~\ref{g26:sec:rigidez-radial-es} fixe, dans
la même réalisation, l'horloge compatible
\[
 t_0=\frac{\pi_{\HMT}L}{54c_{\rm int}}.
\]
Sa représentation circulaire satisfait alors \(R_{108}=L\). Cette
substitution exprime chronologiquement la longueur construite et conserve
la référence longitudinale, l’action et la vitesse de cette réalisation.
Le calendrier déjà construit détermine $T_{108}=108t_0$.
Dans sa réalisation circulaire,
\[
\omega_{108}=\frac{2\pi}{108t_0},\qquad
R_{108}=\frac{c_{\rm int}}{\omega_{108}}
=\frac{54}{\pi}\ell_0,\qquad \ell_0=c_{\rm int}t_0.
\]
L’unité $\ell_0$ et les coordinatisations du temps et de la vitesse sont maintenues
explicites ; la coordinatisation SI ne sélectionne pas le cycle.
Pour chaque section orientée d’action $a\in\{\mathrm{pre},\mathrm{ret}\}$,
\[
E_{108,a}=\hbar_a\omega_{108},\qquad
m_{108,a}=\frac{E_{108,a}}{c_{\rm int}^2}.
\]
Les longueurs réduite et complète sont respectivement
\[
\frac{\hbar_a}{m_{108,a}c_{\rm int}}=R_{108},
\qquad
\frac{h_a}{m_{108,a}c_{\rm int}}=2\pi R_{108}.
\]
La même action et la même durée interviennent dans les deux lectures
\cite{autodual}.

\begin{proposition}[Sélecteur de coïncidence des rayons]
Dans la famille dimensionnelle
\[
G_a(\vartheta)=\vartheta\frac{c_{\rm int}^5t_0^2}{\hbar_a},
\qquad \vartheta>0,
\]
la condition de la réalisation circulaire
$G_a(\vartheta)m_{108,a}/c_{\rm int}^2=R_{108}$
admet l'unique solution
\[
\vartheta^\circ_{108}=(54/\pi)^2.
\]
\end{proposition}
\begin{proof}
Le premier rayon est
$r_g(\vartheta)=\vartheta(\pi/54)\ell_0$.
Le comparer à $(54/\pi)\ell_0$, avec $\ell_0>0$, réduit
l'égalité à une équation linéaire strictement croissante.
\end{proof}
La condition de coïncidence est une partie explicite de cette réalisation ;
le théorème résout son coefficient sans utiliser de valeur métrologique de $G$.
Le rayon employé ici est $Gm/c^2$, de sorte que le facteur deux
de la convention $2Gm/c^2$ n'est pas transféré par inadvertance.

\begin{proposition}[Action et gravité réciproques, aire conservée]
Dans les deux coordinatisations précédentes,
\[
\frac{G_{\rm pre}}{G_{\rm ret}}=R_{\rm act}^{-1},\qquad
\hbar_aG_a=\vartheta^\circ_{108}c_{\rm int}^5t_0^2,
\qquad
\ell_P^2=\frac{\hbar_aG_a}{c_{\rm int}^3}
=\vartheta^\circ_{108}\ell_0^2.
\]
\end{proposition}
\begin{proof}
La définition de $G_a$ contient $\hbar_a^{-1}$ et le reste des facteurs
est commun aux deux coordinatisations. Prendre le quotient et le produit donne
les trois égalités.
\end{proof}
Cette réciprocité explique pourquoi l'information d'orientation peut
varier sans que l'aire change. On dispose ainsi d'une relation précise
entre l'ellipse d'action, l'échelle de Planck et l'opérateur aréal
de la section~\ref{sec:area}.

\subsection{Normalisation thermique et énergie par cycle}
Nous fixons la même section orientée \(a\) pour l’énergie et la
température. La relation thermique du corpus établit
\begin{equation}
k_B\Theta_{{\rm clk},a}\log3
=\frac{h_a}{108t_0}
=\frac{\hbar_a}{t_P}=E_{P,a},
\qquad t_P=\frac{54}{\pi}t_0.
\label{eq:thermalclock}
\end{equation}
Le facteur $\log3$ mesure l’information du triplet équiprobable ;
l’énergie de l’horloge incorpore l’action déjà engendrée.
L’identité conserve le produit $k_B\Theta_{\rm clk}$ comme section
d’énergie et rend explicite la coordinatisation thermique utilisée \cite{termica}.

La réalisation radiative fournit, à cette échelle,
\[
n_\gamma\ell_P^3=\frac{2\zeta(3)}{\pi^2\log^3 3},\qquad
\frac{u_\gamma\ell_P^3}{E_P}
=\frac{\pi^2}{15\log^4 3},\qquad
\frac{s_\gamma\ell_P^3}{k_B}
=\frac{4\pi^2}{45\log^3 3}.
\]
On rencontre ici deux invariants distincts : $\log3$ conserve
l’information TRIT et $\zeta(3)$ le moment d’occupation bosonique.
Leur présence dans la même expression relie les constantes thermiques
aux moments spectraux de l’occupation bosonique, sans identifier
ces grandeurs entre elles.

La section~\ref{sec:ii-kb-aut-desarrollo} développe le transducteur
de Boltzmann, son action sur l'information tritique et le changement de
coordonnées thermiques. L'équation~\eqref{eq:thermalclock} fixe le
produit énergie--température qu'utilise cette construction.
La sélection d'une coordonnée individuelle y est distinguée de
l'identité portant sur le produit. Les formules radiatives précédentes
conservent, dans ce même développement, le contenu d'occupation et
d'entropie de la réalisation bosonique.

\subsection{Barbero--Immirzi dans la réalisation de Cartan--Holst}
Une fois $\gamma_{\HMT}$ construit, la source le réalise comme
coefficient de l’action de Holst. Dans la coordinatisation géométrique déclarée,
soit $e^I$ une co-tétrade non dégénérée, $\omega^{IJ}$ une connexion,
$\Sigma^{IJ}=e^I\wedge e^J$ et
$F^{IJ}=\dd\omega^{IJ}+\omega^I{}_K\wedge\omega^{KJ}$.
Avec $\kappa=8\pi G/c^4$, l’action spécialisée est
\begin{equation}
\begin{aligned}
S[e,\omega,\Psi]
={}&\frac1{2\kappa c}\int
\Sigma_{IJ}\wedge(\star F^{IJ}+\gamma_{\HMT}^{-1}F^{IJ})\\
&-\frac{\Lambda}{\kappa c}\int\operatorname{vol}_e
+S_m[e,\omega,\Psi].
\end{aligned}
\label{eq:holst}
\end{equation}
La convention de courant est
$\delta_\omega S_m=-\tfrac12\int\delta\omega^{IJ}\wedge\sigma_{IJ}$.
Pour cette variation, la connexion est fixée à la frontière,
$\delta\omega^{IJ}|_{\partial M}=0$, ou l'on utilise des variations à
support compact dans l'intérieur. Ces conditions suffisantes annulent
le terme de bord résultant de l'intégration de
$\Sigma_{IJ}\wedge(\star+\gamma_{\HMT}^{-1}I)
D_\omega\delta\omega^{IJ}$.
La variation de \eqref{eq:holst} sous cette condition donne
\[
D_\omega[(\star+\gamma_{\HMT}^{-1}I)\Sigma^{IJ}]
=\kappa c\,\sigma^{IJ}.
\]
On l'obtient en utilisant $\delta F=D_\omega\delta\omega$, en intégrant par parties
et en annulant le coefficient de la variation arbitraire
\cite{holst}. L'orientation de la frontière et la normalisation
du courant font partie de cette formule.

En signature lorentzienne, $\star^2=-I$ sur les bivecteurs. L'identité
\[
(\star+\gamma^{-1}I)^{-1}
=\frac{\gamma^2}{1+\gamma^2}(-\star+\gamma^{-1}I)
\]
se vérifie par multiplication des deux membres : les termes croisés
s'annulent et il reste $(1+\gamma^{-2})I$ avant normalisation.
Ainsi, le paramètre déjà construit intervient dans la réponse torsionnelle,
et non uniquement dans une valeur tabulée.
Les identités
$D_\omega T^I=F^I{}_J\wedge e^J$ et $D_\omega F^{IJ}=0$
maintiennent le lien entre torsion, courbure et bilan.

La réalisation géométrique conserve la mémoire des chemins et la
soudure de frontière. La constante cosmologique \(\Lambda\) de
\eqref{eq:holst} est le coefficient déclaré dans l’action ; elle ne
s’identifie pas à la représentation vectorielle \(\Lambda(R)\).
Les identités suivantes conservent les conditions d’holonomie,
de lissité et de non-dégénérescence de la réalisation.
% Artículo VII, recuperación aditiva, revisión 04, 10-09-2026.
% RESULTADO_RECUPERADO: S15, 11c_gravedad_holonomia_rev6.tex:16--126.
% El cociclo canónico c27=(1,18) se conserva como antecedente del cómputo
% de rutas; este fragmento prueba su composición y representación.
% La memoria traslacional finita y la amplitud cosmológica s0 tienen
% tipos distintos. El cociclo no evalúa por sí solo una densidad local.

\subsection{Retour de phase et représentation de la mémoire accumulée}
\label{vii:sec:retorno-memoria-gravedad}

La chronologie des chemins enrichis distingue le retour de position,
le retour d’orientation et le retour de phase. La projection observable
identifie certains états d’arrivée, tandis que le registre cumulatif
conserve la séquence parcourue. La représentation géométrique doit
transporter conjointement les deux données.

\subsubsection{Cocycle d'inversion et changement de feuillet}
\label{vii:subsec:cociclo-memoria}

Soit \(F_{\rm obs}\) la représentation chronologique du transport TPK.
Ses blocs canoniques ont la hiérarchie
\[
\begin{array}{c|l}
27&\text{retour positionnel avec inversion de l’orientation},\\
54&\text{retour positionnel et orientationnel},\\
108&\text{retour complet de la phase observable}.
\end{array}
\]
Le registre compte les inversions d'orientation \(g(a)\) et les
changements effectifs de feuillet \(s(a)\). Pour une route,
\begin{equation}
 c(\gamma)=\sum_{a\subset\gamma}(g(a),s(a)),\qquad
 c(\gamma_2\gamma_1)=c(\gamma_2)+c(\gamma_1).
 \label{vii:eq:cociclo-ruta}
\end{equation}
Le calcul canonique du bloc de longueur \(27\) donne
\(c_{27}=(1,18)\). La composition conserve ces accroissements dans
les quatre blocs qui complètent la période observable.
Sur les chemins orientés vers l’avant, le registre est cumulatif ;
son extension au groupoïde des chemins utilise \(\mathbb Z^2\) et
\(c(\bar\gamma)=-c(\gamma)\).

\begin{proposition}[Mémoire du tour observable complet]
\label{vii:prop:memoria-vuelta-completa}
En écrivant \(\mathcal K_{27}=F_{\rm obs}^{27}\), le relèvement
au registre, pour \(x\in\mathcal Q_{\rm obs}\) et \(z\in\mathbb Z^2\), satisfait
\[
 \widetilde{\mathcal K}_{27}(x,z)
   =(\mathcal K_{27}x,z+(1,18)),\qquad
 \widetilde{\mathcal K}_{27}^{\,4}(x,z)
   =(x,z+(4,72)).
\]
Après \(N\) tours complets, l'incrément est \(N(4,72)\).
\end{proposition}

\begin{proof}
La projection observable de \(\mathcal K_{27}\) est d'ordre quatre. Le calcul canonique précédent fixe \(c_{27}=(1,18)\) de manière constante sur cette orbite observable. L'additivité de \eqref{vii:eq:cociclo-ruta} additionne les quatre incréments et donne \(c_{108}=4(1,18)=(4,72)\). La répétition applique la même loi de concaténation et produit \(Nc_{108}\). Le retour de la première coordonnée conserve l'incrément de la seconde.
\end{proof}

\subsubsection{Représentation affine du cocycle}
\label{vii:subsec:representacion-afin-cociclo}

La réalisation discrète utilise un élément \(R_4\) d'ordre quatre dans
\(\operatorname{Spin}^+(1,3)\), un vecteur temporel unitaire \(u\)
fixé par son action vectorielle et une unité de longueur \(\ell_0>0\).
L'échelle \(\ell_0\) conserve sa provenance dans la section dimensionnelle.
Nous définissons
\begin{equation}
 \rho_{\rm C}^{\rm disc}(\gamma)=
 \bigl(R_4^{c_g(\gamma)},\,\ell_0c_s(\gamma)u\bigr)
 \in\operatorname{Spin}^+(1,3)\ltimes\mathbb R^{1,3}.
 \label{vii:eq:representacion-discreta-memoria}
\end{equation}
La puissance agit sur le vecteur au moyen de la représentation vectorielle
associée du groupe de spin.

\begin{theorem}[Composition et amplitude translationnelle du retour]
\label{vii:thm:traslacion-memoria}
L'application \(\rho_{\rm C}^{\rm disc}\) conserve l'identité,
la concaténation et l'inversion. Dans les retours complets,
\begin{equation}
 \rho_{\rm C}^{\rm disc}(\gamma_{108})=(I,72\ell_0u),\qquad
 \rho_{\rm C}^{\rm disc}(\gamma_{108}^{\,N})=(I,72N\ell_0u).
 \label{vii:eq:traslacion-retorno-108}
\end{equation}
L'inversion change le signe de la translation.
\end{theorem}

\begin{proof}
Le produit semi-direct est
\((R,a)(S,b)=(RS,a+Rb)\). Comme \(R_4u=u\),
\[
\begin{split}
 \rho_{\rm C}^{\rm disc}(\gamma_2)
 \rho_{\rm C}^{\rm disc}(\gamma_1)
 &=
 \bigl(R_4^{c_g(\gamma_2)+c_g(\gamma_1)},
 \ell_0(c_s(\gamma_2)+c_s(\gamma_1))u\bigr)\\
 &=\rho_{\rm C}^{\rm disc}(\gamma_2\gamma_1).
\end{split}
\]
Le cocycle de l'identité est nul et celui de l'inverse est son opposé.
Pour le tour complet, on substitue \(c_{108}=(4,72)\) et
\(R_4^4=I\). La concaténation de \(N\) tours maintient la partie
rotationnelle égale à l'identité et additionne leurs translations.
\end{proof}

L'amplitude \(72\ell_0\) est une longueur associée à une route finie. Une réalisation au moyen d'une cotétrade et d'une connexion la transporte vers son holonomie affine. La contribution locale aux équations gravitationnelles s'obtient ensuite au moyen de la courbure de cette connexion, de la variation de l'action matérielle et de la normalisation de sa densité. La représentation conserve ainsi le contenu cumulatif exact qui doit intervenir dans cette réalisation.

% Artículo VII. Incorporación de resultados preexistentes, 10-09-2026.
% Propietarios íntegros preservados en fuentes_conservadas:
% 17_pantalla_cubica_soldadura.tex (P), 18_clausura_energetica.tex (E).
% P: incidencia, cociente, soldadura celular y refinamiento.
% E: respuesta, unicidad, extensión mínima y expectativa tracial.
% Dependencias anteriores de VII: núcleo APP--TRIT--TPK; realización
% cúbica de la ruta, cociclo de transporte y compresión q_vis=5/9.
% Este fragmento no constituye por sí solo el artículo autónomo.
% La identidad de pantalla se distingue de la contracción material de espín.

\subsection{Incidence cubique et soudure de frontière}
\label{vii:sec:pantalla-respuesta}

La réalisation cubique du transport conserve six incidences orientées.
Son quotient géométrique, la forme lorentzienne et la réponse énergétique
s’obtiennent à partir de la même matrice d’incidence. Le transport provient du
chemin enrichi APP--TRIT--TPK : la face, le signe, la retenue et la
mémoire accompagnent chaque arête. Cette section développe
l’opération linéaire qui agit sur cette représentation de frontière.

\subsubsection{Incidence cubique et quotient de rang quatre}
\label{vii:sec:cociente-pantalla}

Soit \(H_F=\mathbb R^F\), avec
\(F=\{(i,\epsilon):1\leq i\leq3,\ \epsilon\in\{-1,1\}\}\),
son produit scalaire euclidien et sa base \(e_{i,\epsilon}\).
L'opposition des faces est l'involution orthogonale
\(Re_{i,\epsilon}=e_{i,-\epsilon}\). Nous définissons
\[
 u=\sum_{i,\epsilon}e_{i,\epsilon},\qquad
 P_u=\frac{uu^{\mathsf T}}6,\qquad P_-=\frac{I-R}{2},\qquad
 P_0=\frac{I+R}{2}-P_u,\qquad P_\square=P_u+P_-.
\]
Les sommets sont \(V=\{-1,1\}^3\). La matrice d'incidence \(M:\mathbb R^V\longrightarrow H_F\) est déterminée par
\[
 M_{(i,\epsilon),\nu}=\mathbf1_{\{\nu_i=\epsilon\}}.
\]

\begin{proposition}[Décomposition de l'incidence]
\label{vii:prop:incidencia}
On a
\[
 MM^{\mathsf T}=12P_u+4P_-,
 \qquad
 \ker M^{\mathsf T}=P_0H_F.
\]
Par conséquent, le quotient
\(Q_\square=Q=H_F/P_0H_F\), identifié orthogonalement à
\(P_\square H_F\), est de dimension quatre.
\end{proposition}

\begin{proof}
Une face contient quatre sommets. Deux faces opposées n'ont pas de sommets
communs et deux faces d'axes différents ont deux sommets communs.
En écrivant \(\mathbf J=uu^{\mathsf T}\), on obtient
\[
 MM^{\mathsf T}=4I+2(\mathbf J-I-R)
 =2(I-R)+2\mathbf J=4P_-+12P_u.
\]
Les trois projecteurs \(P_u,P_-,P_0\) sont orthogonaux et leur somme vaut \(I\) ; leurs rangs sont respectivement un, trois et deux. L'identité \(\|M^{\mathsf T}x\|^2=\langle x,MM^{\mathsf T}x\rangle\) identifie le noyau annoncé. Le rang complémentaire est quatre.
\end{proof}

Sur \(Q\), la forme
\[
 B=P_--P_u
\]
a pour signature \((3,1)\). L'orientation temporelle est fixée par \(u\).
Avec
\[
 t=\frac{u}{\sqrt6},\qquad
 d_i=\frac{e_{i,+}-e_{i,-}}{\sqrt2},\qquad
 W=\sqrt6P_u+\sqrt2P_-,
\]
on a \(We_{i,\epsilon}=t+\epsilon d_i\). Ces six vecteurs
sont nuls pour \(B\), tandis que
\[
 3t+\nu_1d_1+\nu_2d_2+\nu_3d_3
\]
a pour norme quadratique \(-6\). La même incidence satisfait
\(MM^{\mathsf T}=2W^*W\) sur \(Q\). Les rotations propres du cube
commutent avec les projecteurs, préservent les deux formes et fixent \(t\).
Ces identités déterminent la métrique de la représentation cubique.

\subsubsection{Soudure cellulaire et transport de mémoire}
\label{vii:sec:soldadura-celular}

Pour chaque arête orientée \(a\), la lecture du chemin fournit
la face \(\lambda(\underline a)\), le signe
\(\sigma(a)\), les fibres d’origine et de destination et le transport
\(U(a):Q_{s(a)}\longrightarrow Q_{t(a)}\).
L’inversion vérifie \(U(\bar a)=U(a)^{-1}\).
Les demi-arêtes de frontière conservent leur inversion formelle et leur incidence
sur le bord. Avec \(\ell_m=\ell_0\,9^{-m}>0\), la soudure est
\begin{equation}
 \beta_m(a)=\ell_m\sigma_m(a)
              n_{\lambda(\underline a)}^{s(a)},\qquad
 n_{i,\epsilon}=t+\epsilon d_i.
 \label{vii:eq:soldadura-arista}
\end{equation}
La face et le signe sont transportés avec la mémoire de la route. La convention d'inversion exige
\begin{equation}
 \beta_m(\bar a)=-U_m(a)\beta_m(a).
 \label{vii:eq:inversion-soldadura}
\end{equation}
En appliquant deux fois cette relation, on retrouve \(\beta_m(a)\) ; le registre de retenue conserve séparément la composition ordonnée.

\begin{proposition}[Naturalité de la soudure par raffinement]
\label{vii:prop:refinamiento-soldadura}
Soit \(a=a_1\cdots a_9\) un raffinement compatible d’une arête.
Supposons que les neuf incidences filles, transportées vers la fibre du
prédécesseur, conservent la face et le signe, et que
\(\ell_{m+1}=\ell_m/9\). Soit
\(\mathcal U_j:Q_{m,s(a)}\longrightarrow Q_{m+1,s(a_j)}\)
le transport cumulé qui inclut l’identification par raffinement
et le trajet jusqu’à l’origine du successeur \(j\). Alors
\[
 \beta_m(a)=
 \sum_{j=1}^{9}\mathcal U_j^{-1}\beta_{m+1}(a_j).
\]
L'égalité est compatible avec l'inversion et avec la concaténation du registre de mémoire.
\end{proposition}

\begin{proof}
Chaque terme, exprimé dans la fibre de l’origine du prédécesseur, est
\((\ell_m/9)\sigma_m(a)n_{\lambda(\underline a)}\).
La somme de neuf termes reproduit la soudure du prédécesseur.
L’inversion renverse l’ordre des transports et remplace chacun
par son inverse ; la relation \eqref{vii:eq:inversion-soldadura} fournit
le signe global. La mémoire se concatène dans ce même ordre, de sorte
que le raffinement conserve tant la représentation vectorielle que son
registre enrichi.
\end{proof}

% RESULTADO_RECUPERADO: S15, 11c_gravedad_holonomia_rev6.tex:128--212;
% S01, 02_dinamica_tres_hojas.tex:155--219; pantalla y respuesta de30.
% Pruebas de curvatura, covariancia y Bianchi reunidas para autonomía.
% La realización suave conserva sus condiciones de holonomía explícitas.

\subsection{Réalisation géométrique du transport et de la soudure}
\label{vii:sec:realizacion-cartan}

La représentation discrète du retour et la soudure d'écran déterminent deux aspects du transport : l'accumulation de mémoire dans les routes et la comparaison de leurs fibres. Leur réalisation géométrique conserve les deux opérations. Dans cette section sont fixés l'espace d'arrivée, la normalisation dimensionnelle de la connexion et les identités qu'utilisera sa variation.

\subsubsection{Écran, cotétrade et connexion}

Le quotient d'écran \(Q_\square\), sa forme lorentzienne et la
soudure \(\beta_m\) procèdent de
\ref{vii:sec:pantalla-respuesta}. Dans chaque cellule, une réalisation d'
écran \(J_{{\rm scr},m}\) transporte la forme et la soudure
vers la fibre géométrique correspondante. Son image cellulaire est
\[
 e_m=J_{{\rm scr},m}\beta_m.
\]
L'échelle \(\ell_m\) demeure à l'intérieur de la soudure. Les
changements de repère agissent sur \(J_{{\rm scr},m}\), la connexion
et le transport des cellules conjointement.

Une réalisation lisse de ces données sur une variété orientée \(\mathcal U\) de dimension quatre consiste en une cotétrade non dégénérée \(e=(e^I)\), une connexion de Lorentz \(\omega=(\omega^I{}_J)\) et une réalisation de routes \(\mathfrak R_{\rm C}\) qui satisfont
\begin{equation}
 \operatorname{Hol}_{\mathcal A_{\ell_0}}
       (\mathfrak R_{\rm C}(\gamma))
 =\rho_{\rm C}(\operatorname{Hol}_{\rm TPK}(\gamma)).
 \label{vii:eq:compatibilidad-holonomia-geometrica}
\end{equation}
Dans cette égalité matricielle, la représentation est normalisée par
\[
 N_{\ell_0}(R,a)=(\Lambda(R),a/\ell_0),\qquad
 \rho_{\rm C}=N_{\ell_0}\circ\rho_{\rm C}^{\rm disc},
\]
où \(\Lambda:\operatorname{Spin}^+(3,1)\to SO^+(3,1)\)
est la représentation vectorielle. L’application \(N_{\ell_0}\)
préserve le produit semi-direct : diviser la translation
\(a+\Lambda(R)b\) par \(\ell_0\) équivaut à composer les
deux translations normalisées. Par conséquent, le tour complet produit
\(72u\) dans la connexion normalisée et \(72\ell_0u\) dans la
coordinatisation dimensionnelle. Le relèvement de spin conserve séparément
\(R_4^{c_g(\gamma)}\) ; la représentation vectorielle a pour
noyau central \(\{1,-1\}\), de sorte que ce
relèvement est retenu lorsqu’on utilise des champs spinoriels.
L’identité, la concaténation, l’inversion et la subdivision sont maintenues dans la
réalisation correspondante. L’appariement de
l’écran et la condition d’holonomie spécifient les données de la
réalisation utilisée dans les équations de champ.

Avec \(\eta=\operatorname{diag}(-1,1,1,1)\), on définit
\(g=\eta_{IJ}e^I\otimes e^J\). Pour une échelle élémentaire
positive \(\ell_0\), constante dans la coordinatisation, la connexion affine est
\begin{equation}
 \mathcal A_{\ell_0}
 =\begin{pmatrix}\omega&\ell_0^{-1}e\\0&0\end{pmatrix}.
 \label{vii:eq:conexion-afin-realizada}
\end{equation}
La cotétrade a la dimension d'une longueur ; le facteur \(\ell_0^{-1}\) homogénéise sa composante translationnelle.

\begin{proposition}[Courbure et torsion de la connexion réalisée]
\label{vii:prop:curvatura-afin}
La courbure de \eqref{vii:eq:conexion-afin-realizada} est
\begin{equation}
 \mathcal F_{\ell_0}
 =\mathrm d\mathcal A_{\ell_0}
   +\mathcal A_{\ell_0}\wedge\mathcal A_{\ell_0}
 =\begin{pmatrix}F_\omega&\ell_0^{-1}T\\0&0\end{pmatrix},
 \quad F_\omega=\mathrm d\omega+\omega\wedge\omega,
 \quad T=\mathrm de+\omega\wedge e.
 \label{vii:eq:curvatura-afin-realizada}
\end{equation}
\end{proposition}
\begin{proof}
Le produit de matrices de formes produit \(\omega\wedge\omega\) dans le bloc diagonal et \(\ell_0^{-1}\omega\wedge e\) dans le bloc translationnel. La dérivée extérieure apporte \(\mathrm d\omega\) et \(\ell_0^{-1}\mathrm de\), puisque \(\ell_0\) est constant. La somme démontre l'égalité.
\end{proof}

La courbure décrit le défaut de transport des repères et la torsion
celui de la cotétrade. Toutes deux appartiennent à la même connexion réalisée ;
leurs composantes conservent des domaines tensoriels distincts.

\begin{proposition}[Covariance et identités de Bianchi]
\label{vii:prop:bianchi-cartan}
Sous un changement de repère \(g_L:\mathcal U\to SO^+(3,1)\),
\[
 e'=g_Le,\qquad
 \omega'=g_L\omega g_L^{-1}-\mathrm dg_L\,g_L^{-1},
\]
on a \(T'=g_LT\) et \(F_{\omega'}=g_LF_\omega g_L^{-1}\). En outre,
\begin{equation}
 D_\omega T=F_\omega\wedge e,\qquad D_\omega F_\omega=0.
 \label{vii:eq:bianchi-realizacion}
\end{equation}
\end{proposition}
\begin{proof}
Dans \(\mathrm d(g_Le)+\omega'\wedge g_Le\), les termes
contenant \(\mathrm dg_L\wedge e\) se compensent.
La même substitution dans la courbure donne sa transformation par
conjugaison. Pour la première identité de Bianchi, on développe
\[
 D_\omega(\mathrm de+\omega\wedge e)
 =(\mathrm d\omega+\omega\wedge\omega)\wedge e.
\]
Dans \(\mathrm dF_\omega+\omega\wedge F_\omega-F_\omega\wedge\omega\), les termes contenant \(\mathrm d\omega\) et les produits cubiques s'annulent par paires. Cela prouve la seconde identité.
\end{proof}

\subsubsection{Les trois réalisations quadratiques de courbure}

Le feuillet \(\tau\in\{+1,0,-1\}\) du TRIT est représenté dans
l'algèbre de Cartan au moyen des générateurs \(J_{ab}=-J_{ba}\)
et \(P_a^{(\tau)}\), avec les relations
\begin{align*}
 [J_{ab},J_{cd}]&=
 \eta_{bc}J_{ad}-\eta_{ac}J_{bd}
 -\eta_{bd}J_{ac}+\eta_{ad}J_{bc},\\
 [J_{ab},P_c^{(\tau)}]&=
 \eta_{bc}P_a^{(\tau)}-\eta_{ac}P_b^{(\tau)},\\
 [P_a^{(\tau)},P_b^{(\tau)}]&=-\tau J_{ab}.
\end{align*}
Pour \(\ell>0\) constant, on écrit
\[
 \mathcal A_\tau^{\rm Cart}
 =\tfrac12\omega^{ab}J_{ab}+\ell^{-1}e^aP_a^{(\tau)}.
\]

\begin{proposition}[Courbure de la réalisation tritifiée]
\label{vii:prop:curvatura-tritificada}
Avec les relations précédentes,
\begin{equation}
 \mathcal F_\tau^{\rm Cart}
 =\tfrac12(F_\omega^{ab}-\tau\ell^{-2}e^a\wedge e^b)J_{ab}
     +\ell^{-1}T^aP_a^{(\tau)}.
 \label{vii:eq:curvatura-tritificada}
\end{equation}
\end{proposition}
\begin{proof}
On développe \(\mathrm d\mathcal A+\tfrac12[\mathcal A,\mathcal A]\). Le crochet du secteur de Lorentz produit \(F_\omega\), le crochet mixte produit \(D_\omega e\), et le crochet translationnel apporte \(-\tfrac12\tau\ell^{-2}e^a\wedge e^bJ_{ab}\).
\end{proof}

Dans le feuillet central, on retrouve la connexion affine ; les feuillets extrêmes
conservent les deux signes du terme de courbure constante. La
représentation maintient ainsi le régime tritique de la route avant
d'appliquer les équations variationnelles. L'égalité d'holonomie
\eqref{vii:eq:compatibilidad-holonomia-geometrica} correspond au
feuillet central. Dans chaque feuillet extrême, on utilise la représentation
\(\rho_{{\rm C},\tau}\) dans son groupe de Cartan et sa condition
d'holonomie respective. Le courant et le tenseur métrique utilisent
les champs et la représentation de la même réalisation.


\subsection{Transport orienté de la connexion et de l'aire}
\label{sec:ii-costura-barbero}
La parité de la fonctionnelle d’incidence permet de comparer les coordinatisations
\(\gamma\) et \(-\gamma\). Dans la réalisation canonique géométrique,
soit \(E^a_i\) une triade densitisée, \(\Gamma_a^i(E)\) sa
connexion intrinsèque et \(K_a^i\) la courbure extrinsèque dans la
convention normale choisie. La connexion d’Ashtekar--Barbero est
\(\mathcal A_a^i=\Gamma_a^i(E)+\gamma K_a^i\).

\begin{proposition}[Involution entre coordinatisations orientées]
\label{prop:ii-costura-barbero}
Pour \(\gamma\in\mathbb R\setminus\{0\}\), l'application
\[
 \mathfrak b:(\gamma,K,E)\longmapsto(-\gamma,-K,E)
\]
est involutive. Elle conserve \(\mathcal A\) et inverse la forme
\[
 \Omega_{KE}=\frac1\kappa\int_\Sigma
 \delta K_a^i\wedge\delta E^a_i,
\]
avec \(\kappa\) fixe et les mêmes conditions de frontière.
\end{proposition}
\begin{proof}
Appliquer deux fois \(\mathfrak b\) restitue les trois coordonnées.
Puisque \(E\) reste fixe, \(\Gamma(E)\) reste également fixe.
Le produit des deux coordonnées impaires satisfait
\((-\gamma)(-K)=\gamma K\) et conserve la connexion.
Dans l'espace des variations, \(\delta(-K)=-\delta K\) et
\(\delta E\) reste fixe : le tiré en arrière de la forme est
\(\mathfrak b^*\Omega_{KE}=-\Omega_{KE}\).
\end{proof}

L’aire orientée \(\mathcal A_{\rm or}=\gamma a_0N\) et sa
grandeur physique \(\mathcal A_{\rm fis}=|\gamma|a_0N\), pour
\(a_0>0\) et \(N\geq0\), coïncident dans la coordinatisation positive.
Lors du prolongement à la coordinatisation opposée, la première est impaire et la seconde
demeure invariante. Cette distinction conserve tant l’orientation
de l’incidence que la positivité de l’observable géométrique.
La proposition établit le transport algébrique de ces coordinatisations ;
sa réalisation comme inversion de la normale spatiale utilise,
en outre, le transport de la matière et des données de frontière
de l’action~\eqref{eq:holst}.
